\section{Open research question: expansion and unequal string relaxation}
\label{sec:v500-string-relaxation}

The following comparison records the open question at the v5 release
boundary. It is an analytical exploration, not a native consequence of
the graded carrier or a physical identification of fundamental strings.
Suppose three extensible inner paths have lengths $\ell_i$, while an
outer path has fixed length $L$. The equality $\sum_i\ell_i=L$ can be
imposed throughout motion or only as a terminal condition; these are
different models. A concrete network, its winding classes and its
endpoint constraints have not yet been specified.

\paragraph{Elastic allocation and its limitation.}
For independent nonnegative lengths, minimizing
$U=\tfrac12\sum_i k_i\ell_i^2$, with $k_i>0$ and fixed sum $L$,
gives $\ell_i/L=k_i^{-1}/\sum_j k_j^{-1}$ and equal generalized
tensions. Under the additional readout $\sqrt{m_i}=c\ell_i$, $c>0$,
the Koide ratio is $K=\sum_i\ell_i^2/L^2$.
Equal stiffnesses give $K=1/3$. Stiffnesses $(k/q,k,k)$ give
$K=(q^2+2)/(q+2)^2$, so $q=4+3\sqrt2$ yields $K=2/3$.
This factor is inferred from the target, not independently predicted;
the example also has two equal masses. Moreover, $K>1/2$ requires one
length greater than $L/2$. The paths therefore cannot be the three
sides of a straight triangle or three pairwise shortest geodesics in
one metric. Prescribed winding, guides or another network would be needed.

\paragraph{A length scale associated with a cosmological constant.}
Dimensional consistency permits $\Lambda=C/L^2$, but does not determine
$C$ or a stress--energy tensor. If the outer loop is additionally
identified with a great circle on a pure four-dimensional de Sitter
horizon, its radius $R_\Lambda=\sqrt{3/\Lambda}$ gives
$L=2\pi R_\Lambda$ and $C=12\pi^2$.
This fixed horizon scale is not a material loop comoving with expansion.
If instead $L=Aa(t)$ and a separately conserved homogeneous,
isotropic component has $\rho=B/L^2$, with constants $A,B>0$, the continuity equation
$\dot\rho+3H(\rho+p)=0$ gives $w=p/\rho=-1/3$, in units $c=1$.
Vacuum energy requires $w=-1$. Three individual paths do not by
themselves supply the required isotropic cosmological component.

\paragraph{A dissipative comparison.}
Let $r_i$ be an evolving unloaded length. In the tensile regime, a
Maxwell-type model with constant $k_i,\eta_i>0$ is
\begin{equation}
 F_i=k_i(\ell_i-r_i),\qquad
 \dot r_i=F_i/\eta_i,\qquad
 \dot F_i=k_i\dot\ell_i-F_i/\tau_i,\quad \tau_i=\eta_i/k_i.
\end{equation}
Different relaxation times allow different responses to the same imposed
extension. Tension can decrease while length increases, provided
$k_i\dot\ell_i<F_i/\tau_i$. At fixed length,
$F_i(t)=F_i(0)e^{-t/\tau_i}$. The elastic energy satisfies
\begin{equation}
 \dot U_i=F_i\dot\ell_i-F_i^2/\eta_i,\qquad
 U_i=\tfrac12 k_i(\ell_i-r_i)^2.
\end{equation}
The dissipated energy must enter the full energy accounting. Slack paths
require a unilateral tensile law; negative spring force must not be
interpreted as compressive string tension. This comparison represents
creep and relaxation, not a derived cyclic-fatigue law or microscopic
damage mechanism.

If $\sum_i\ell_i=L$ holds at every instant, then
$\sum_i\dot\ell_i=0$: all three actual lengths cannot grow together,
although their unloaded lengths may grow and their tensions may relax.
For freely redistributed quasistatic lengths, force balance can equalize
tensions even when material responses differ. A soft penalty
$\kappa(L-\sum_i\ell_i)^2/2$ supplies the same generalized force to
each length, and does not impose exact equality of the lengths' sum.

\paragraph{The retained question.}
Can a justified coupling between expansion, internal relaxation and a
fixed outer length select a stable unequal allocation? Answering it
requires an embedded network, an expansion-to-strain map, an energy
balance including dissipative degrees of freedom, and a mechanism that
limits relaxation. Identical initial data and identical deterministic
laws alone do not select distinct responses. No value of $\Lambda$,
lepton spectrum, or dynamically selected $K=2/3$ follows here. All
existing physical-identification guards remain open.

\paragraph{Sources for the comparison.}
The de Sitter radius is discussed by H.-Y.~Guo,
\emph{The Beltrami Model of De Sitter Space} (2006),
arXiv:hep-th/0607017. For vacuum cosmology see S.~M.~Carroll,
\emph{A No-Nonsense Introduction to General Relativity} (2001),
Section~6. For finite-rate viscoelastic comparisons see C.-Y.~Lin
et al., \emph{Polymers} \textbf{14} (2022), 2124,
doi:10.3390/polym14102124. These sources support the background models,
not their proposed identification with the Companion's physical sectors.
